\section{4. Electrical validation of the wireless stimulator}\label{electrical-validation-of-the-wireless-stimulator}

\subsection{4.1 Current accuracy and pulse fidelity}\label{current-accuracy-and-pulse-fidelity}

Both outputs of a device running firmware v0.5 were loaded with 1 kΩ resistors, and one output was recorded with a Joulescope JS220 Precision Energy Analyzer at 500 kHz for $\sim$11 s at amplitude settings of 10, 50 and 100 \% (180 µs, 130 Hz; Figure 5). The same three settings were then recorded with both outputs loaded with nominal 47 kΩ resistors, within the range of effective loads estimated in vivo. Delivered current was taken as the voltage across the load divided by the load resistance. The meter's own current channel showed 4--18 µs spikes after each edge from automatic range switching, so both channels were passed through a 38 µs median filter, which removes impulses while preserving steps; charge computed from the voltage changed by less than 0.5 nanocoulombs (nC) with filtering.

Delivered amplitude increased with the setting and was highly repeatable from pulse to pulse. The stimulating phase measured 106, 264 and 530 µA at 10, 50 and 100 \%, with a pulse-to-pulse SD below 0.5 µA at every setting (Figure 5B). A line through the three settings has a slope of 4.7 µA per percent and an intercept of +47 µA, against the 6 µA per percent shown in the application. The application value is therefore a reproducible setting rather than a calibrated current: the device delivers more than the nominal value at low settings and less at high settings. The recovery phase scaled with the stimulating phase at 1:4.7--4.9 (22, 56 and 113 µA).

Pulse timing differed from the set values by a similar amount at every amplitude. At the 180 µs setting the stimulating phase lasted 238--254 µs and the recovery phase 854--878 µs, against 900 µs set, so the two phases did not carry equal charge: the median pulse delivered 27, 63 and 125 nC in the stimulating phase and 19, 48 and 96 nC in the recovery phase, leaving a net 9.5, 17.3 and 33.0 nC per period in the stimulating direction (26--35 \% of the stimulating-phase charge; Figure 5C). The 130 Hz setting delivered 142.5 Hz at every amplitude, a 7.02 ms period, because firmware v0.5 sets the period in whole milliseconds (Table 2), and the period varied with an SD of 29--38 µs, about one 30.5 µs timer tick.

Into 47 kΩ, delivered current was limited by compliance (Figure 5B). At the 50 and 100 \% settings the stimulating phase was flat-topped at 92.6 µA, which implies an output swing of 4.6 V across the load and the 2.66 kΩ of sense and filter resistance in series with it. The recovery phase reached 47 µA at 50 \%, below the 56 µA delivered into 1 kΩ, and was also limited at 100 \%, at about 88 µA (4.4 V). At 10 \%, below the limit, the edges were slow at this load: the stimulating phase was still rising when it ended (peak 51 µA, against 106 µA into 1 kΩ), and the cause, possibly capacitance on the load side, was not identified. Because these effects reduce the stimulating phase more than the recovery phase, the net charge per period changed sign with amplitude, at +5.5, $-$13.3 and $-$48.4 nC for 10, 50 and 100 \% (Figure 5C). At a 0 \% setting into 47 kΩ the output carried a residual of about 1 nC per period during the active window, small relative to the imbalance above; in use, stimulation is turned off rather than set to 0 \%. Other pulse widths and loads were not characterized.

\subsection{4.2 Battery and operating duration}\label{battery-and-operating-duration}

Battery current was measured at the cell terminals in each operating state with the same Joulescope JS220 (Table 3). The device draws 1.47 mA when idle and not advertising, 1.62 mA while advertising, 1.89 mA while connected to the application, and 3.15--3.42 mA while stimulating at amplitude settings labeled 60 and 400 µA, respectively, into 1 kΩ with a 180 µs pulse width at the 130 Hz frequency setting (142.5 Hz delivered).

\textbf{Table 3.} Measured battery current and projected runtime on two size-13 zinc-air cells in series ($\sim$280 mAh; 90 \% of nameplate capacity assumed usable). Stimulation rows use a 1 kΩ load, a 180 µs pulse width and the 130 Hz setting, which delivers 142.5 Hz, and give the amplitude as labeled in the application. The continuous-stimulation projection was confirmed in a manual bench test that was not logged.

\begin{longtable}{@{}>{\raggedright\arraybackslash{}}p{0.4\linewidth}>{\raggedleft\arraybackslash{}}p{0.22\linewidth}>{\raggedleft\arraybackslash{}}p{0.25\linewidth}@{}}
\toprule
Operating state & Battery current & Projected runtime \\
\midrule
\endhead
Idle, not advertising & 1.47 mA & $\sim$7.1 d \\
Advertising & 1.62 mA & $\sim$6.5 d \\
Connected to app & 1.89 mA & $\sim$5.6 d \\
Stimulating, 60 µA & 3.15 mA & $\sim$3.3 d \\
Stimulating, 400 µA & 3.42 mA & $\sim$3.1 d \\
\bottomrule
\end{longtable}

Two features of these numbers define the usable experimental window. First, raising the amplitude setting from 60 to 400 µA increased battery current by only 0.27 mA, so delivered amplitude contributes relatively little to the power budget. Entering the stimulation state adds 1.68--1.95 mA over idle; this reflects circuitry and processing that are active while the stimulation protocol runs, and how it divides among subsystems was not measured. Second, even at the highest measured current a fresh pair of cells has a projected runtime of about three days of continuous stimulation, and more than five days of standby with the phone connected. The power budget therefore permits hours-long, overnight and multiday sessions without disturbing the animal, with longer studies possible through a brief battery exchange, which requires detaching the device and pauses stimulation; charge balance in the present firmware is a separate constraint. Zinc-air cells are inexpensive and widely available, which suits a device whose cells are replaced frequently.

The device is intended for chronic studies with scheduled battery exchange, not as a fully implanted, maintenance-free system.
